\section{Implementación del Algoritmo de Criba}
Una vez definido el cambio de paradigma hacia un modelo de pre-cálculo, se procedió a desarrollar la solución en Python. A continuación, se presentan las dos etapas de optimización alcanzadas.

\subsection{Versión en Python Nativo}
Esta primera implementación traduce la lógica de la criba de sumas utilizando las estructuras de datos estándar del lenguaje. La clave de su rendimiento reside en el uso del tercer parámetro de la función $range$, que permite iterar únicamente sobre los múltiplos de cada divisor.

\begin{lstlisting}[language=Python]
def amigos(MAX):
    t1 = time.time()

    # Inicializacion: El 0 no tiene divisores (0).
    # Se precarga 1 para el resto (divisor universal).
    suma_divisores = [0] + [1] * MAX

    # Criba de divisores: Reparto de cada 'divisor' entre sus multiplos
    for divisor in range(2, MAX // 2 + 1):
        # El salto 'divisor' evita el uso del operador modulo (%)
        for multiplo in range(divisor * 2, MAX + 1, divisor):
            suma_divisores[multiplo] += divisor

    # Verificacion de pares: Identifica perfectos y pares simetricos
    for numero in range(MAX + 1):
        posible_amigo = suma_divisores[numero]
        
        # Validacion de reciprocidad: 
        # Si la suma de divisores del 'posible_amigo' es igual al
        #'numero' original.
        if posible_amigo <= MAX and suma_divisores[posible_amigo] == numero:
            print(numero, posible_amigo)

    print(time.time() - t1)
\end{lstlisting}

A diferencia del modelo original, donde el trabajo dependía de evaluar cada número individualmente ($n \times n$), la eficiencia de la criba reside en que el número de operaciones disminuye drásticamente a medida que el divisor aumenta. El análisis se divide en las tres etapas del código:

\begin{itemize}
    \item \textbf{Inicialización de la Estructura:}
    La instrucción \colorbox{gray!10}{\texttt{[0, 0] + [1] * (MAX - 1)}} crea una lista de tamaño $n$ en un solo paso de memoria. Esta operación tiene una complejidad lineal:
    \[ T_{init}(n) = O(n) \]

    \item \textbf{Cuerpo del Algoritmo (Reparto de Divisores):}
    Esta es la parte central. El bucle externo recorre los divisores $d$ desde $2$ hasta $n/2$. Por cada divisor, el bucle interno realiza una cantidad de sumas igual a la cantidad de múltiplos de $d$ que entran en el rango MAX.
    \begin{itemize}
        \item Para $d=2$, se realizan $n/2$ operaciones.
        \item Para $d=3$, se realizan $n/3$ operaciones.
        \item Para $d=4$, se realizan $n/4$ operaciones.
        \item ... hasta $d=n/2$, donde se realizan $2$ operaciones.
    \end{itemize}
    La suma total de operaciones es:
    \[ 
    T(n) = n \times \left( \frac{1}{2} + \frac{1}{3} + \frac{1}{4} + \dots + \frac{1}{n/2} \right) 
    \]
    Esta expresión representa una variante de la \textbf{Serie Armónica}. Según el análisis de sumatorias presentado por \textbf{Cormen et al. (2022)}, la suma de los recíprocos de los primeros n términos se comporta de forma logarítmica:
    \[ \sum_{k=1}^{n} \frac{1}{k} \approx \ln n + \gamma \]
    Por lo tanto, la complejidad dominante de esta etapa, y por extensión del algoritmo completo, es:
    \[
    T_{criba}(n) = O(n \log n)
    \]

    \item \textbf{Verificación de Pares:}
    El último bucle recorre la lista de sumas ya calculada una sola vez para comparar los valores y encontrar las reciprocidades. Al ser un recorrido simple:
    \[ T_{verif}(n) = O(n) \]
\end{itemize}

\textbf{Complejidad Final:} Sumando las etapas y despreciando los términos de menor orden, se concluye que la complejidad total del algoritmo optimizado es:
\[ \boxed{\textbf{Complejidad} = \bm{O(n \log n)}} \]

\subsection{Optimización Avanzada con NumPy}
Aunque la versión nativa reduce la complejidad a $O(n\log n)$, el intérprete de Python sigue encontrando un cuello de botella al gestionar los bucles for internos. Para alcanzar un rendimiento de grado industrial, se implementó una versión utilizando la librería NumPy.
Esta biblioteca permite delegar el trabajo pesado a rutinas pre-compiladas en C, utilizando una técnica denominada vectorización.

\begin{lstlisting}[language=Python]

import time
import numpy as np


def amigos(MAX):
    t1 = time.time()

    tabla_sumas = np.ones(MAX + 1, dtype=np.int64)
    tabla_sumas[0] = 0
    tabla_sumas[1] = 0

    for divisor in range(2, MAX // 2 + 1):
        tabla_sumas[divisor * 2 :: divisor] += divisor

    numeros = np.arange(MAX + 1)
    es_candidato = (tabla_sumas > numeros) & (tabla_sumas <= MAX)
    lista_candidatos = np.where(es_candidato)[0]

    for n in lista_candidatos:
        suma_n = tabla_sumas[n]
        if tabla_sumas[suma_n] == n:
            print(n, suma_n)

    print(time.time() - t1)
\end{lstlisting}

\begin{itemize}
    \item \textbf{Slicing Vectorizado:} La instrucción \colorbox{gray!10}{\lstinline[language=Python]|tabla_sumas[divisor * 2 :: divisor] += divisor|} elimina el overhead del bucle \texttt{for} de Python. NumPy delega la iteración a rutinas pre-compiladas en C, que acceden directamente a bloques de memoria contiguos y realizan las sumas en paralelo mediante instrucciones \textbf{SIMD} (Single Instruction, Multiple Data).
    
    \item \textbf{Indexación Booleana:} En lugar de recorrer los $n$ elementos buscando candidatos
    con un \texttt{if} tradicional, se utiliza la expresión: \\
    \colorbox{gray!10}{\lstinline[language=Python]|(tabla_sumas > numeros) & (tabla_sumas <=
    MAX)|}.\\
    Esto crea una ``máscara'' de bits que filtra el arreglo completo a nivel de hardware, reduciendo drásticamente la cantidad de comparaciones que el intérprete de Python debe procesar.
    
    \item \textbf{Filtrado Lógico de Resultados:} Mediante la condición \colorbox{gray!10}{\lstinline[language=Python]|tabla_sumas < numeros|}, el algoritmo realiza dos limpiezas críticas en un solo paso:
        \begin{itemize}
            \item \textbf{Eliminación de Perfectos:} Al exigir que la suma sea menor (o distinta) al número ($s < n$), se excluyen casos como el 6 o el 28, donde $s = n$.
            \item \textbf{Garantía de Unicidad:} Al procesar solo los casos donde el ``amigo'' es menor al número actual, se evita informar el mismo par dos veces (ej. muestra $[220, 284]$ pero no $[284, 220]$), reduciendo la redundancia en la salida final.
        \end{itemize}
\end{itemize}

\textbf{Nota sobre la complejidad:} Aunque el uso de NumPy no cambia la complejidad asintótica teórica ($O(n \log n)$), la constante de tiempo se reduce de forma tan significativa que permite procesar en milisegundos lo que a Python nativo le tomaría segundos.
